\documentclass[11pt,a4paper]{article}
\usepackage[utf8]{inputenc}
\usepackage[vietnamese]{babel}
\usepackage{amsmath,amssymb}
\usepackage[version=4]{mhchem}
\usepackage{geometry}
\geometry{top=1.5cm,bottom=1.2cm,left=1.4cm,right=1.4cm,headheight=28pt,headsep=12pt,footskip=18pt}
\usepackage{graphicx}
\usepackage{xcolor}
\usepackage{tabularx}
\usepackage{array}
\usepackage{tikz}
\usepackage{fancyhdr}
\usepackage{enumitem}
\usepackage{multido}
\usepackage{tcolorbox}

\definecolor{hfRed}{HTML}{D32F2F}
\definecolor{hfGreen}{HTML}{8BC390}
\definecolor{mainGreen}{HTML}{2E7D32}
\definecolor{solBlue}{HTML}{1565C0}
\definecolor{tableGreen}{HTML}{E8F5E9}

\usepackage{ex_test}

\pagestyle{fancy}
\fancyhf{}
\renewcommand{\headrulewidth}{0.75pt}
\renewcommand{\headrule}{\hbox to\headwidth{\color{hfGreen}\leaders\hrule height \headrulewidth\hfill}}
\renewcommand{\footrulewidth}{0.75pt}
\renewcommand{\footrule}{\hbox to\headwidth{\color{hfGreen}\leaders\hrule height \footrulewidth\hfill}}

\fancyhead[L]{\small\itshape\color{hfRed}Tài liệu lưu hành nội bộ}
\fancyhead[R]{\small\color{hfRed}Thầy \textbf{TRẦN VĂN THẠNH} - 0777.470.803}

\fancyfoot[L]{\footnotesize\itshape\color{hfRed}CS1: 6/15 Nguyễn Hoàng, P. Kim Long, TP Huế \quad|\quad CS2: 24 Đặng Thái Thân, TP Huế}
\fancyfoot[R]{\small\textbf{\color{hfRed}Trang \thepage}}

\setlength{\parindent}{0pt}
\setlength{\parskip}{2pt}

\begin{document}

\begin{minipage}[t]{0.48\textwidth}
\textbf{SỞ GD\&ĐT THỪA THIÊN HUẾ}\\
\textbf{TRƯỜNG THPT NGUYỄN ĐÌNH CHIỂU}\\
\textit{ĐỀ CHÍNH THỨC}\\
\textit{(Đề thi có 02 trang)}
\end{minipage}
\hfill
\begin{minipage}[t]{0.48\textwidth}
\raggedleft
\textbf{\color{hfRed}KIỂM TRA GIỮA HỌC KỲ I}\\
\textbf{MÔN: HOÁ HỌC -- LỚP 10}\\
\textit{Thời gian làm bài: 45 phút (không kể thời gian giao đề)}\\
\textbf{MÃ ĐỀ: ĐỀ GỐC}
\end{minipage}

\vspace{4pt}
\noindent
\textit{Họ, tên thí sinh:} \dotfill \textit{Lớp:} \dotfill \textit{Số báo danh:} \dotfill

\vspace{3pt}
\noindent
\textit{\footnotesize Cho nguyên tử khối: $\ce{H} = 1; \ce{C} = 12; \ce{N} = 14; \ce{O} = 16; \ce{Na} = 23; \ce{Mg} = 24; \ce{Al} = 27; \ce{P} = 31; \ce{S} = 32; \ce{Cl} = 35,5; \ce{K} = 39; \ce{Cr} = 52; \ce{Fe} = 55,85; \ce{Cu} = 64; \ce{Br} = 80; \ce{Ba} = 137$. Số Avogadro $N_A = 6,022 \times 10^{23}\text{ mol}^{-1}$.}

\vspace{3pt}
\hrule color{hfGreen} height 0.5pt

\vspace{4pt}
\begin{center}
\textbf{\color{hfRed}\large (BẢN GIÁO VIÊN -- CÓ LỜI GIẢI CHI TIẾT)}
\end{center}

\vspace{2pt}
\begin{center}
{\color{hfRed}\textbf{\large BẢNG ĐÁP ÁN RÚT GỌN (DÀNH CHO GIÁO VIÊN)}}
\end{center}

\vspace{2pt}
\noindent
\textbf{\color{mainGreen}1. Phần I: Trắc nghiệm nhiều lựa chọn (3,0 điểm - mỗi câu đúng 0,25 điểm)}
\begin{center}
\renewcommand{\arraystretch}{1.2}
\begin{tabular}{|c|c|c|c|c|c|c|c|c|c|c|c|c|}
\hline
\rowcolor{tableGreen}
\textbf{Câu} & \textbf{1} & \textbf{2} & \textbf{3} & \textbf{4} & \textbf{5} & \textbf{6} & \textbf{7} & \textbf{8} & \textbf{9} & \textbf{10} & \textbf{11} & \textbf{12} \\
\hline
\textbf{Đ/A} & \textbf{\color{hfRed}A} & \textbf{\color{hfRed}B} & \textbf{\color{hfRed}A} & \textbf{\color{hfRed}A} & \textbf{\color{hfRed}C} & \textbf{\color{hfRed}D} & \textbf{\color{hfRed}B} & \textbf{\color{hfRed}B} & \textbf{\color{hfRed}D} & \textbf{\color{hfRed}C} & \textbf{\color{hfRed}A} & \textbf{\color{hfRed}A} \\
\hline
\end{tabular}
\end{center}

\vspace{2pt}
\noindent
\textbf{\color{mainGreen}2. Phần II: Trắc nghiệm Đúng / Sai (3,0 điểm)}
\begin{center}
\renewcommand{\arraystretch}{1.2}
\begin{tabular}{|c|l|c|c|c|c|}
\hline
\rowcolor{tableGreen}
\textbf{Câu} & \textbf{Nội dung câu hỏi} & \textbf{Lệnh a} & \textbf{Lệnh b} & \textbf{Lệnh c} & \textbf{Lệnh d} \\
\hline
\textbf{Câu 1} & Đồng vị $\ce{Cl}$ và $\ce{Cu}$ & \textbf{\color{hfRed}ĐÚNG} & \textbf{\color{hfRed}ĐÚNG} & \textbf{\color{hfRed}ĐÚNG} & \textbf{SAI} \\
\hline
\textbf{Câu 2} & Các nguyên tố chu kì 3 & \textbf{\color{hfRed}ĐÚNG} & \textbf{SAI} & \textbf{\color{hfRed}ĐÚNG} & \textbf{\color{hfRed}ĐÚNG} \\
\hline
\textbf{Câu 3} & Cấu hình e và ô orbital & \textbf{\color{hfRed}ĐÚNG} & \textbf{\color{hfRed}ĐÚNG} & \textbf{SAI} & \textbf{\color{hfRed}ĐÚNG} \\
\hline
\end{tabular}
\end{center}

\vspace{2pt}
\noindent
\textbf{\color{mainGreen}3. Phần III: Trả lời ngắn (1,0 điểm - mỗi câu đúng 0,25 điểm)}
\begin{center}
\renewcommand{\arraystretch}{1.2}
\begin{tabular}{|c|c|c|c|c|}
\hline
\rowcolor{tableGreen}
\textbf{Câu} & \textbf{Câu 1} & \textbf{Câu 2} & \textbf{Câu 3} & \textbf{Câu 4} \\
\hline
\textbf{Đáp án} & \textbf{\color{hfRed}137} & \textbf{\color{hfRed}124} & \textbf{\color{hfRed}11} & \textbf{\color{hfRed}2} \\
\hline
\end{tabular}
\end{center}

\vspace{2pt}
\noindent
\textbf{\color{mainGreen}4. Tóm tắt đáp số Phần Tự luận (3,0 điểm)}
\begin{center}
\renewcommand{\arraystretch}{1.2}
\begin{tabular}{|c|c|p{11.5cm}|}
\hline
\rowcolor{tableGreen}
\textbf{Câu} & \textbf{Điểm} & \textbf{Đáp số cốt lõi} \\
\hline
\textbf{Câu 1} & 1,0 đ & a) $p = e = 12$; $n = 12$. \quad b) Cấu hình: $\ce{1s^2 2s^2 2p^6 3s^2}$; là kim loại. \\
\hline
\textbf{Câu 2} & 1,0 đ & a) $\ce{X}$ ($Z = 11$, CK 3, nhóm IA); $\ce{Y}$ ($Z = 12$, CK 3, nhóm IIA). \quad b) Base: $\ce{NaOH} > \ce{Mg(OH)2}$. \\
\hline
\textbf{Câu 3} & 1,0 đ & a) Tên: Sulfur (Lưu huỳnh), KHHH: $\ce{S}$. \quad b) $m = 9,8\text{ gam acid } (\ce{H2SO4})$. \\
\hline
\end{tabular}
\end{center}

\vspace{4pt}
\hrule color{hfGreen} height 0.5pt

\vspace{4pt}
\noindent
{\color{mainGreen}\textbf{\Large A. PHẦN TRẮC NGHIỆM (7 điểm)}}

\vspace{3pt}
\noindent
{\color{mainGreen}\textbf{\large PHẦN I (3 điểm). CÂU TRẮC NGHIỆM NHIỀU PHƯƠNG ÁN LỰA CHỌN}}\\
\textit{\small Thí sinh trả lời từ câu 1 đến câu 12. Mỗi câu hỏi thí sinh chỉ chọn 1 phương án.}
\vspace{2pt}

\begin{ex}
Nguyên tử của nguyên tố $\ce{X}$ có tổng số hạt cơ bản (p, n, e) là 52, trong đó số hạt mang điện nhiều hơn số hạt không mang điện là 16. Số khối của nguyên tử $\ce{X}$ là
\choice
{\True 35}
{36}
{34}
{37}

\loigiai{
Gọi $p, n, e$ lần lượt là số proton, neutron, electron của $\ce{X}$ ($p = e = Z$).\\
Ta có hệ phương trình:
$\begin{cases} 2Z + N = 52 \\ 2Z - N = 16 \end{cases} \implies 4Z = 68 \implies Z = 17, N = 18$.\\
Số khối của $\ce{X}$: $A = Z + N = 17 + 18 = 35$.\\
\textbf{Chọn A.}
}
\end{ex}
\begin{ex}
Hợp chất khí với hydrogen của nguyên tố $\ce{R}$ có dạng $\ce{RH3}$. Trong oxide cao nhất của $\ce{R}$, oxygen chiếm 74,07\% về khối lượng. Phát biểu nào sau đây sai?
\choice
{Nguyên tố $\ce{R}$ thuộc chu kì 2, nhóm VA}
{\True Phân tử oxide cao nhất của $\ce{R}$ có công thức $\ce{R2O3}$}
{Độ âm điện của $\ce{R}$ lớn hơn của phosphorus}
{Hydroxide ứng với hóa trị cao nhất của $\ce{R}$ là $\ce{HNO3}$}

\loigiai{
Hợp chất khí với hydrogen có dạng $\ce{RH3} \implies \ce{R}$ thuộc nhóm VA.\\
Oxide cao nhất của $\ce{R}$ có dạng $\ce{R2O5}$.\\
$\%O = \dfrac{16 \times 5}{2M_{\ce{R}} + 16 \times 5} \times 100\% = 74,07\% \implies \dfrac{80}{2M_{\ce{R}} + 80} = 0,7407 \implies M_{\ce{R}} = 14 \text{ (Nitrogen - N)}$.\\
- A đúng: $\ce{N}$ ($Z = 7$) thuộc chu kì 2, nhóm VA.\\
- B sai: Công thức oxide cao nhất là $\ce{R2O5}$ ($\ce{N2O5}$), không phải $\ce{R2O3}$.\\
- C đúng: $\ce{N}$ và $\ce{P}$ cùng nhóm VA, chu kì nhỏ hơn có độ âm điện lớn hơn $\implies$ độ âm điện của $\ce{N} > \ce{P}$.\\
- D đúng: Hydroxide cao nhất là $\ce{HNO3}$.\\
\textbf{Chọn B.}
}
\end{ex}
\begin{ex}
Nguyên tử của nguyên tố chromium ($\ce{Cr}$) có số hiệu nguyên tử $Z = 24$. Cấu hình electron ở trạng thái cơ bản của nguyên tử $\ce{Cr}$ là
\choice
{\True $\ce{[Ar] 3d^5 4s^1}$}
{$\ce{[Ar] 3d^5 4s^2}$}
{$\ce{[Ar] 3d^4 4s^2}$}
{$\ce{[Ar] 3d^6}$}

\loigiai{
$\ce{Cr}$ ($Z = 24$): Theo trật tự năng lượng là $\ce{1s^2 2s^2 2p^6 3s^2 3p^6 4s^2 3d^4}$.\\
Do hiện tượng bán bão hòa phân lớp d bền vững hơn, 1 electron chuyển từ $4s$ sang $3d$, cấu hình electron trạng thái cơ bản là: $\ce{1s^2 2s^2 2p^6 3s^2 3p^6 3d^5 4s^1}$ hay $\ce{[Ar] 3d^5 4s^1}$.\\
\textbf{Chọn A.}
}
\end{ex}
\begin{ex}
Ở trạng thái cơ bản, phân lớp 3d chứa tối đa bao nhiêu electron?
\choice
{\True 10}
{6}
{18}
{14}

\loigiai{
Phân lớp d gồm 5 orbital (AO), mỗi orbital chứa tối đa 2 electron.\\
Do đó, phân lớp 3d chứa tối đa $5 \times 2 = 10$ electron.\\
\textbf{Chọn A.}
}
\end{ex}
\begin{ex}
Cho các nguyên tố: $\ce{C} (Z = 6), \ce{N} (Z = 7), \ce{Si} (Z = 14), \ce{P} (Z = 15)$. Chiều bán kính nguyên tử tăng dần là
\choice
{$\ce{Si < P < C < N}$}
{$\ce{N < P < C < Si}$}
{\True $\ce{N < C < P < Si}$}
{$\ce{C < N < Si < P}$}

\loigiai{
- Chu kì 2: $\ce{C}$ (nhóm IVA), $\ce{N}$ (nhóm VA). Trong chu kì: bán kính $\ce{N < C}$.\\
- Chu kì 3: $\ce{Si}$ (nhóm IVA), $\ce{P}$ (nhóm VA). Trong chu kì: bán kính $\ce{P < Si}$.\\
- Trong nhóm: bán kính chu kì 2 nhỏ hơn chu kì 3: $\ce{C < Si}$ và $\ce{N < P}$.\\
$\implies$ Thứ tự tăng dần bán kính: $\ce{N < C < P < Si}$.\\
\textbf{Chọn C.}
}
\end{ex}
\begin{ex}
Cho các nguyên tố $\ce{Na} (Z = 11), \ce{Mg} (Z = 12), \ce{Al} (Z = 13), \ce{K} (Z = 19)$. Chiều tính base của các hydroxide tăng dần là
\choice
{$\ce{Al(OH)3 < Mg(OH)2 < KOH < NaOH}$}
{$\ce{Mg(OH)2 < Al(OH)3 < NaOH < KOH}$}
{$\ce{KOH < NaOH < Mg(OH)2 < Al(OH)3}$}
{\True $\ce{Al(OH)3 < Mg(OH)2 < NaOH < KOH}$}

\loigiai{
- Trong chu kì 3: tính kim loại $\ce{Al < Mg < Na} \implies$ tính base: $\ce{Al(OH)3 < Mg(OH)2 < NaOH}$.\\
- Trong nhóm IA: $\ce{K}$ ở chu kì 4 có tính kim loại mạnh hơn $\ce{Na} \implies \ce{KOH > NaOH}$.\\
$\implies$ Chiều tính base tăng dần: $\ce{Al(OH)3 < Mg(OH)2 < NaOH < KOH}$.\\
\textbf{Chọn D.}
}
\end{ex}
\begin{ex}
Trong tự nhiên, nguyên tố bromine ($\ce{Br}$) có hai đồng vị bền là $^{79}\ce{Br}$ và $^{81}\ce{Br}$. Biết nguyên tử khối trung bình của bromine là 79,91. Phần trăm số nguyên tử của đồng vị $^{79}\ce{Br}$ là
\choice
{49,5\%}
{\True 54,5\%}
{50,5\%}
{45,5\%}

\loigiai{
Gọi $x$ là phần trăm của đồng vị $^{79}\ce{Br}$, suy ra $^{81}\ce{Br}$ chiếm $(100 - x)\%$.\\
Ta có: $\overline{A} = \dfrac{79x + 81(100 - x)}{100} = 79,91 \implies 8100 - 2x = 7991 \implies 2x = 109 \implies x = 54,5\%$.\\
\textbf{Chọn B.}
}
\end{ex}
\begin{ex}
Theo chiều tăng dần của điện tích hạt nhân trong một chu kì, xu hướng nào sau đây là đúng đối với các nguyên tố nhóm A?
\choice
{Độ âm điện giảm dần}
{\True Tính kim loại giảm dần, tính phi kim tăng dần}
{Bán kính nguyên tử tăng dần}
{Tính base tăng dần}

\loigiai{
Trong một chu kì, theo chiều tăng của điện tích hạt nhân (từ trái sang phải):\\
- Bán kính nguyên tử giảm dần;\\
- Độ âm điện tăng dần;\\
- Tính kim loại giảm dần, tính phi kim tăng dần;\\
- Tính base của oxide/hydroxide giảm dần, tính acid tăng dần.\\
\textbf{Chọn B.}
}
\end{ex}
\begin{ex}
Nguyên tử của nguyên tố $\ce{Y}$ ở trạng thái cơ bản có tổng số electron ở các phân lớp p là 11. Vị trí của $\ce{Y}$ trong bảng tuần hoàn là
\choice
{chu kì 3, nhóm VIA}
{chu kì 4, nhóm IA}
{chu kì 3, nhóm VA}
{\True chu kì 3, nhóm VIIA}

\loigiai{
Tổng số electron ở các phân lớp p là 11 $\implies$ gồm $\ce{2p^6}$ và $\ce{3p^5}$.\\
Cấu hình electron đầy đủ của $\ce{Y}$: $\ce{1s^2 2s^2 2p^6 3s^2 3p^5}$ ($Z = 17$).\\
- Có 3 lớp electron $\implies$ Chu kì 3.\\
- Lớp ngoài cùng có 7 electron ($2 + 5 = 7$), nguyên tố p $\implies$ Nhóm VIIA.\\
\textbf{Chọn D.}
}
\end{ex}
\begin{ex}
Hạt nhân của hầu hết các nguyên tử được cấu tạo từ các loại hạt nào sau đây?
\choice
{Electron và neutron}
{Proton và electron}
{\True Proton và neutron}
{Proton, neutron và electron}

\loigiai{
Hạt nhân của hầu hết các nguyên tử được cấu tạo từ hai loại hạt cơ bản: proton (mang điện tích dương) và neutron (không mang điện). Electron chuyển động ở lớp vỏ.\\
\textbf{Chọn C.}
}
\end{ex}
\begin{ex}
Trong bảng tuần hoàn các nguyên tố hóa học, số thứ tự của nhóm A bằng
\choice
{\True số electron ở lớp ngoài cùng}
{tổng số hạt proton và neutron}
{số electron ở phân lớp ngoài cùng}
{số lớp electron trong nguyên tử}

\loigiai{
Đối với các nguyên tố nhóm A, số thứ tự của nhóm bằng số electron ở lớp ngoài cùng.\\
\textbf{Chọn A.}
}
\end{ex}
\begin{ex}
Cặp kí hiệu nguyên tử nào sau đây là đồng vị của cùng một nguyên tố hóa học?
\choice
{\True $^{35}_{17}\ce{X}$ và $^{37}_{17}\ce{Y}$}
{$^{40}_{19}\ce{X}$ và $^{40}_{20}\ce{Y}$}
{$^{16}_{8}\ce{X}$ và $^{14}_{7}\ce{Y}$}
{$^{23}_{11}\ce{X}$ và $^{24}_{12}\ce{Y}$}

\loigiai{
Đồng vị là các nguyên tử có cùng số proton (cùng $Z$) nhưng khác số neutron (khác $A$).\\
Cặp $^{35}_{17}\ce{X}$ và $^{37}_{17}\ce{Y}$ có cùng $Z = 17$, số khối khác nhau (35 và 37) nên là đồng vị của cùng một nguyên tố (Chlorine).\\
\textbf{Chọn A.}
}
\end{ex}
\vspace{4pt}
\noindent
{\color{mainGreen}\textbf{\large PHẦN II (3 điểm). CÂU TRẮC NGHIỆM ĐÚNG / SAI}}\\
\textit{\small Thí sinh trả lời từ câu 1 đến câu 3. Trong mỗi ý a), b), c), d) ở mỗi câu, thí sinh chọn đúng hoặc sai.}
\vspace{2pt}

\begin{ex}
Trong tự nhiên, nguyên tố chlorine ($\ce{Cl}$) có hai đồng vị bền là $^{35}\ce{Cl}$ (75,77\%) và $^{37}\ce{Cl}$ (24,23\%). Nguyên tố copper ($\ce{Cu}$) có hai đồng vị bền là $^{63}\ce{Cu}$ và $^{65}\ce{Cu}$ với nguyên tử khối trung bình là 63,54. (Khi tính khối lượng mol của $\ce{CuCl2}$, dùng nguyên tử khối trung bình của $\ce{Cl}$ tính từ phần trăm các đồng vị ở trên.)
\choiceTF
{\True Hạt nhân nguyên tử đồng vị $^{37}\ce{Cl}$ có 17 proton và 20 neutron}
{\True Trong hợp chất $\ce{CuCl2}$, phần trăm khối lượng của đồng vị $^{63}\ce{Cu}$ xấp xỉ bằng 34,19\%}
{\True Nguyên tử khối trung bình của chlorine là 35,4846}
{Nếu lấy 13,45092 gam $\ce{CuCl2}$ thì tổng số hạt neutron trong các nguyên tử $^{37}\ce{Cl}$ là $2,91826 \times 10^{23}$ hạt}

\loigiai{
a) \textbf{Đúng}: Đồng vị $^{37}_{17}\ce{Cl}$ có $Z = 17$ (17 proton) và $N = 37 - 17 = 20$ neutron.\\
b) \textbf{Đúng}: Gọi $x$ là \% số nguyên tử $^{63}\ce{Cu}$: $\overline{A}_{\ce{Cu}} = \dfrac{63x + 65(100 - x)}{100} = 63,54 \implies x = 73\%$.\\
Khối lượng mol của $\ce{CuCl2} = 63,54 + 2 \times 35,4846 = 134,5092\text{ g/mol}$.\\
Trong 1 mol $\ce{CuCl2}$ có 0,73 mol $^{63}\ce{Cu} \implies m_{^{63}\ce{Cu}} = 0,73 \times 63 = 45,99\text{ g}$.\\
$\%m_{^{63}\ce{Cu}} = \dfrac{45,99}{134,5092} \times 100\% \approx 34,19\%$.\\
c) \textbf{Đúng}: $\overline{A}_{\ce{Cl}} = \dfrac{35 \times 75,77 + 37 \times 24,23}{100} = 35,4846$.\\
d) \textbf{Sai}: $n_{\ce{CuCl2}} = \dfrac{13,45092}{134,5092} = 0,1\text{ mol} \implies n_{\ce{Cl}} = 0,1 \times 2 = 0,2\text{ mol}$.\\
$n_{^{37}\ce{Cl}} = 0,2 \times 24,23\% = 0,04846\text{ mol}$.\\
Số hạt neutron trong $^{37}\ce{Cl} = 0,04846 \times 20 \times 6,022 \times 10^{23} \approx 5,8365 \times 10^{23}\text{ hạt}$ (khác $2,91826 \times 10^{23}$).
}
\end{ex}
\begin{ex}
Cho các nguyên tố thuộc chu kì 3 trong bảng tuần hoàn: $\ce{Na, Mg, Al, P, S, Cl}$.
\choiceTF
{\True Trong số các nguyên tố trên, $\ce{Na}$ là nguyên tố có bán kính nguyên tử lớn nhất và $\ce{Cl}$ có độ âm điện lớn nhất}
{Các oxide cao nhất gồm $\ce{Na2O, MgO, Al2O3}$ đều là các basic oxide, còn $\ce{P2O5, SO3, Cl2O7}$ đều là các acidic oxide}
{\True Theo chiều tăng dần của điện tích hạt nhân, hóa trị cao nhất của các nguyên tố với oxygen tăng dần từ I đến VII, đồng thời hóa trị với hydrogen của $\ce{P, S, Cl}$ giảm dần từ III xuống I}
{\True Cho một lượng dư dung dịch $\ce{NaOH}$ vào hỗn hợp gồm $\ce{Al(OH)3}$ và $\ce{Mg(OH)2}$, sau phản ứng chỉ có $\ce{Al(OH)3}$ bị hòa tan}

\loigiai{
a) \textbf{Đúng}: Đi từ trái sang phải trong chu kì 3, bán kính giảm dần ($\ce{Na}$ lớn nhất) và độ âm điện tăng dần ($\ce{Cl}$ lớn nhất).\\
b) \textbf{Sai}: $\ce{Al2O3}$ là oxide lưỡng tính (amphoteric oxide), không phải basic oxide.\\
c) \textbf{Đúng}: Hóa trị cao nhất với oxygen tăng từ I đến VII ($\ce{Na2O}$ đến $\ce{Cl2O7}$). Hóa trị với hydrogen của $\ce{P, S, Cl}$ lần lượt là III ($\ce{PH3}$), II ($\ce{H2S}$), I ($\ce{HCl}$) giảm dần từ III xuống I.\\
d) \textbf{Đúng}: $\ce{Al(OH)3}$ là hydroxide lưỡng tính phản ứng tan trong kiềm: $\ce{Al(OH)3 + NaOH -> NaAlO2 + 2H2O}$. $\ce{Mg(OH)2}$ là base không tan trong $\ce{NaOH}$.
}
\end{ex}
\begin{ex}
Cho bốn nguyên tố $\ce{X, Y, Z, T}$ có số hiệu nguyên tử lần lượt là $7, 14, 24, 26$.
\choiceTF
{\True Khi biểu diễn cấu hình electron theo ô orbital ở trạng thái cơ bản, nguyên tử $\ce{T} (Z = 26)$ có tổng cộng 15 orbital chứa electron, trong đó có 4 AO chứa 1 electron độc thân}
{\True Nguyên tử $\ce{Y}$ ở trạng thái cơ bản có 4 electron ở lớp ngoài cùng và thuộc chu kì 3, nhóm IVA}
{Ở trạng thái cơ bản, nguyên tử $\ce{Z} (Z = 24)$ có 6 electron độc thân và có 2 electron ở lớp ngoài cùng}
{\True Ở trạng thái cơ bản, nguyên tử $\ce{X}$ có 3 electron độc thân ở phân lớp 2p}

\loigiai{
a) \textbf{Đúng}: $\ce{T} (Z = 26, \ce{Fe})$: $\ce{1s^2 2s^2 2p^6 3s^2 3p^6 3d^6 4s^2}$. Số AO chứa e = $1 + 1 + 3 + 1 + 3 + 5 + 1 = 15$ AO. Phân lớp $\ce{3d^6}$ có 1 cặp ghép đôi và 4 electron độc thân.\\
b) \textbf{Đúng}: $\ce{Y} (Z = 14, \ce{Si})$: $\ce{1s^2 2s^2 2p^6 3s^2 3p^2}$, có 3 lớp e $\implies$ chu kì 3, có 4e lớp ngoài cùng $\implies$ nhóm IVA.\\
c) \textbf{Sai}: $\ce{Z} (Z = 24, \ce{Cr})$: $\ce{[Ar] 3d^5 4s^1}$ có 6e độc thân, nhưng lớp ngoài cùng (lớp 4) chỉ có 1 electron ($\ce{4s^1}$), không phải 2 electron.\\
d) \textbf{Đúng}: $\ce{X} (Z = 7, \ce{N})$: $\ce{1s^2 2s^2 2p^3}$, 3 electron ở $\ce{2p}$ phân bố vào 3 AO nên có 3 electron độc thân.
}
\end{ex}
\vspace{4pt}
\noindent
{\color{mainGreen}\textbf{\large PHẦN III (1 điểm). CÂU TRẮC NGHIỆM YÊU CẦU TRẢ LỜI NGẮN}}\\
\textit{\small Thí sinh trả lời từ câu 1 đến câu 4.}
\vspace{2pt}

\begin{ex}
Hòa tan hoàn toàn 6,85 gam một kim loại nhóm IIA ($\ce{M}$) vào 200 gam dung dịch $\ce{HCl}$ 3,65\%, cô cạn dung dịch thu được 10,4 gam chất rắn khan. Xác định nguyên tử khối của $\ce{M}$.

\loigiai{
Phương trình phản ứng: $\ce{M + 2HCl -> MCl2 + H2}$\\
$n_{\ce{HCl}} = \dfrac{200 \times 3,65\%}{36,5} = 0,2\text{ mol}$.\\
Cô cạn dung dịch thu được muối khan $\ce{MCl2}$:\\
$m_{\ce{MCl2}} = n_{\ce{M}} \times (M_{\ce{M}} + 71) = 6,85 + 71 \times n_{\ce{M}} = 10,4\text{ g} \implies n_{\ce{M}} = \dfrac{10,4 - 6,85}{71} = 0,05\text{ mol}$.\\
$M_{\ce{M}} = \dfrac{6,85}{0,05} = 137\text{ (Barium - Ba)}$.\\
\textbf{Đáp số: 137}
}
\end{ex}
\begin{ex}
Ở 20~$^\circ\text{C}$, khối lượng riêng của $\ce{Fe}$ ($M = 55,85\text{ g/mol}$) là $7,87\text{ g/cm}^3$. Giả thiết nguyên tử $\ce{Fe}$ là hình cầu chiếm 68\% thể tích tinh thể. Tính bán kính nguyên tử $\ce{Fe}$ theo đơn vị pm (làm tròn kết quả đến hàng đơn vị).

\loigiai{
Thể tích thực của 1 mol nguyên tử $\ce{Fe}$:\\
$V_{\text{thực}} = \dfrac{55,85}{7,87} \times 68\% \approx 4,82567\text{ cm}^3$.\\
Thể tích 1 nguyên tử $\ce{Fe}$:\\
$V_1 = \dfrac{4,82567}{6,022 \times 10^{23}} \approx 8,0134 \times 10^{-24}\text{ cm}^3$.\\
$r = \sqrt[3]{\dfrac{3V_1}{4\pi}} = \sqrt[3]{\dfrac{3 \times 8,0134 \times 10^{-24}}{4 \times 3,14159}} \approx 1,2413 \times 10^{-8}\text{ cm} = 124,13\text{ pm} \approx 124\text{ pm}$.\\
\textbf{Đáp số: 124}
}
\end{ex}
\begin{ex}
Trong tự nhiên, $\ce{H}$ có 2 đồng vị bền ($^1\ce{H}, ^2\ce{H}$), $\ce{S}$ có 1 đồng vị bền ($^{32}\ce{S}$), $\ce{O}$ có 3 đồng vị bền ($^{16}\ce{O}, ^{17}\ce{O}, ^{18}\ce{O}$). Có tối đa bao nhiêu loại phân tử $\ce{H2SO4}$ khác nhau về phân tử khối được tạo thành từ các đồng vị trên?

\loigiai{
$M_{\ce{H2SO4}} = m(\ce{H2}) + 32 + m(\ce{O4})$.\\
- $m(\ce{H2})$ nhận các giá trị: $1+1=2, 1+2=3, 2+2=4 \implies \{2, 3, 4\}$.\\
- $m(\ce{O4})$ nhận các giá trị nguyên liên tiếp từ $16 \times 4 = 64$ đến $18 \times 4 = 72$ (9 giá trị).\\
$\implies M_{\ce{H2SO4}}$ quét qua toàn bộ các số nguyên liên tiếp từ $2 + 32 + 64 = 98$ đến $4 + 32 + 72 = 108$.\\
Số loại phân tử khối khác nhau = $108 - 98 + 1 = 11$ loại.\\
\textbf{Đáp số: 11}
}
\end{ex}
\begin{ex}
Hai nguyên tố $\ce{X, Y}$ thuộc nhóm A, kế tiếp nhau trong một chu kì, tổng số proton trong hai hạt nhân là 31. Ở trạng thái cơ bản, nguyên tử có độ âm điện lớn hơn có bao nhiêu electron độc thân?

\loigiai{
$Z_{\ce{X}} + Z_{\ce{Y}} = 31$ và $|Z_{\ce{Y}} - Z_{\ce{X}}| = 1 \implies Z_{\ce{X}} = 15\ (\ce{P}), Z_{\ce{Y}} = 16\ (\ce{S})$.\\
Trong chu kì 3, độ âm điện $\ce{S} > \ce{P}$.\\
Cấu hình e của $\ce{S} (Z = 16)$: $\ce{1s^2 2s^2 2p^6 3s^2 3p^4}$.\\
Phân lớp $\ce{3p^4}$ có 3 AO: 1 AO ghép đôi và 2 AO chứa electron độc thân $\implies$ Có 2 electron độc thân.\\
\textbf{Đáp số: 2}
}
\end{ex}
\vspace{6pt}
\noindent
{\color{mainGreen}\textbf{\Large B. PHẦN TỰ LUẬN (3,0 điểm)}}
\vspace{2pt}

\noindent
\textbf{Câu 1: (1,0 điểm)} Nguyên tử của nguyên tố $\ce{X}$ có tổng số các loại hạt cơ bản (proton, neutron, electron) là 36. Trong đó, số hạt mang điện gấp đôi số hạt không mang điện.\\
a) Tính số hạt proton, neutron và electron của nguyên tử $\ce{X}$.\\
b) Viết cấu hình electron của nguyên tử $\ce{X}$ và cho biết $\ce{X}$ là kim loại, phi kim hay khí hiếm? Vì sao?

\begin{tcolorbox}[colback=white,colframe=solBlue,title=\textbf{Lời giải chi tiết và Biểu điểm Câu 1:}]
\textbf{a) (0,5 điểm)}
\begin{itemize}[leftmargin=1.2cm]
\item Gọi $p, n, e$ là số hạt của $\ce{X}$ ($p = e$).
\item Hệ phương trình: $\begin{cases} 2p + n = 36 \\ 2p = 2n \iff p = n \end{cases} \hfill \textit{(0,25 đ)}$
\item Giải hệ: $3n = 36 \implies n = 12 \implies p = e = 12$.
\item Vậy số hạt proton là 12, electron là 12, neutron là 12. \hfill \textit{(0,25 đ)}
\end{itemize}
\textbf{b) (0,5 điểm)}
\begin{itemize}[leftmargin=1.2cm]
\item Cấu hình electron của $\ce{X} (Z = 12)$: $\ce{1s^2 2s^2 2p^6 3s^2}$ (hoặc $\ce{[Ne] 3s^2}$). \hfill \textit{(0,25 đ)}
\item Nguyên tử $\ce{X}$ có 2 electron ở lớp ngoài cùng nên $\ce{X}$ là kim loại (Magnesium - Mg). \hfill \textit{(0,25 đ)}
\end{itemize}
\end{tcolorbox}

\vspace{4pt}
\noindent
\textbf{Câu 2: (1,0 điểm)} Cho hai nguyên tố $\ce{X}$ có số hiệu nguyên tử $Z = 11$ (Sodium) và $\ce{Y}$ có số hiệu nguyên tử $Z = 12$ (Magnesium).\\
a) Viết cấu hình electron của $\ce{X, Y}$ và xác định vị trí của chúng (chu kì, nhóm) trong bảng tuần hoàn.\\
b) Viết công thức hydroxide tương ứng của $\ce{X, Y}$. Theo quy luật tuần hoàn, hãy so sánh tính base của hai hydroxide này.

\begin{tcolorbox}[colback=white,colframe=solBlue,title=\textbf{Lời giải chi tiết và Biểu điểm Câu 2:}]
\textbf{a) (0,5 điểm)}
\begin{itemize}[leftmargin=1.2cm]
\item Nguyên tố $\ce{X} (Z = 11)$: Cấu hình $\ce{1s^2 2s^2 2p^6 3s^1}$. Vị trí: Ô 11, chu kì 3, nhóm IA. \hfill \textit{(0,25 đ)}
\item Nguyên tố $\ce{Y} (Z = 12)$: Cấu hình $\ce{1s^2 2s^2 2p^6 3s^2}$. Vị trí: Ô 12, chu kì 3, nhóm IIA. \hfill \textit{(0,25 đ)}
\end{itemize}
\textbf{b) (0,5 điểm)}
\begin{itemize}[leftmargin=1.2cm]
\item Công thức hydroxide tương ứng: của $\ce{X}$ là $\ce{NaOH}$; của $\ce{Y}$ là $\ce{Mg(OH)2}$. \hfill \textit{(0,25 đ)}
\item So sánh: $\ce{Na}$ và $\ce{Mg}$ cùng thuộc chu kì 3, từ trái sang phải theo chiều tăng $Z$, tính kim loại giảm dần $\implies$ tính base giảm dần: $\ce{NaOH > Mg(OH)2}$. \hfill \textit{(0,25 đ)}
\end{itemize}
\end{tcolorbox}

\vspace{4pt}
\noindent
\textbf{Câu 3: (1,0 điểm)} Nguyên tố $\ce{R}$ thuộc nhóm VIA trong bảng tuần hoàn. Trong phân tử oxide cao nhất của $\ce{R}$, nguyên tố oxygen chiếm 60\% về khối lượng.\\
a) Xác định tên và kí hiệu hóa học của nguyên tố $\ce{R}$.\\
b) Hòa tan hoàn toàn 8,0 gam oxide cao nhất của $\ce{R}$ vào nước dư, thu được dung dịch chứa $m$ gam acid. Tính giá trị của $m$.

\begin{tcolorbox}[colback=white,colframe=solBlue,title=\textbf{Lời giải chi tiết và Biểu điểm Câu 3:}]
\textbf{a) (0,5 điểm)}
\begin{itemize}[leftmargin=1.2cm]
\item $\ce{R}$ thuộc nhóm VIA $\implies$ oxide cao nhất là $\ce{RO3}$. \hfill \textit{(0,25 đ)}
\item $\%O = \dfrac{3 \times 16}{M_{\ce{R}} + 48} \times 100\% = 60\% \implies M_{\ce{R}} + 48 = \dfrac{48}{0,6} = 80 \implies M_{\ce{R}} = 32\text{ g/mol}$.
\item Vậy nguyên tố $\ce{R}$ là Sulfur (Lưu huỳnh), kí hiệu hóa học là $\ce{S}$. \hfill \textit{(0,25 đ)}
\end{itemize}
\textbf{b) (0,5 điểm)}
\begin{itemize}[leftmargin=1.2cm]
\item Oxide cao nhất là $\ce{SO3}$ ($M = 80\text{ g/mol}$). Số mol $\ce{SO3}$: $n_{\ce{SO3}} = \dfrac{8,0}{80} = 0,1\text{ mol}$. \hfill \textit{(0,15 đ)}
\item PTHH: $\ce{SO3 + H2O -> H2SO4}$. \hfill \textit{(0,15 đ)}
\item Theo phương trình: $n_{\ce{H2SO4}} = n_{\ce{SO3}} = 0,1\text{ mol}$.
\item Khối lượng acid $\ce{H2SO4}$: $m = 0,1 \times 98 = 9,8\text{ gam}$. Vậy $m = 9,8$. \hfill \textit{(0,2 đ)}
\end{itemize}
\end{tcolorbox}

\vspace{6pt}
\begin{center}
\textbf{\large ---------- HẾT ----------}
\end{center}

\end{document}
